% !TeX root = ../../Book3.tex
\section{\texorpdfstring{\(L^2\)}{L²} 边界值与 Hardy 核}
\label{b3:sec:2802}
径向平方平均的有界性已经等价于系数平方可和。现在把系数放回圆周：每个内部圆周上的 \(f(r\zeta)\) 都含有衰减因子 \(r^n\)，让 \(r\uparrow1\) 就是在逐渐撤去这种衰减。平方可和性保证这个过程在 \(L^2\) 范数中有极限。所得边界函数将把内部点值表示为一个连续线性泛函，核函数由此出现。

\subsection{径向极限与边界函数}
\label{b3:subsec:280201}

对 \(f\in H^2(\mathbb D)\)，把圆周函数 \(\zeta\mapsto f(r\zeta)\) 简记为 \(f_r\)。不同的 \(r\) 对应不同的内部圆周，但参数都放在同一个测度空间 \((\mathbb T,m)\) 上，这样才可以谈论它们的 \(L^2\) 距离。

\begin{theorem}[Hardy 函数的平方积分边界值]\label{b3:thm:28-hardy-boundary}
设 \(f\in H^2(\mathbb D)\)，其 Taylor 系数为 \((a_n)_{n\ge0}\)，圆周测度为 \(\mathrm dm(e^{\mathrm i\theta})=\mathrm d\theta/(2\pi)\)。存在唯一的 \(f^*\in L^2(\mathbb T,m;\mathbb C)\)，使
\begin{equation}\label{b3:eq:28-hardy-radial-limit}
 \lim_{r\uparrow1}\|f_r-f^*\|_{L^2(m)}=0.
\end{equation}
这里的唯一性按几乎处处相等理解。具体地，
\[
 f^*=\lim_{N\to\infty}\sum_{n=0}^Na_n\zeta^n
 \quad\text{于 }L^2(\mathbb T,m),
\]
并且
\begin{equation}\label{b3:eq:28-boundary-error}
 \|f_r-f^*\|_{L^2(m)}^2
 =\sum_{n=0}^{\infty}|a_n|^2(1-r^n)^2,
 \quad \|f^*\|_{L^2(m)}=\|f\|_{H^2}.
\end{equation}
对任意 \(f,g\in H^2(\mathbb D)\)，还有
\[
 \langle f,g\rangle_{H^2}
 =\int_{\mathbb T}f^*(\zeta)\cdot\overline{g^*(\zeta)}\,\mathrm dm(\zeta).
\]
称 \(f^*\) 为 \(f\) 的 \(L^2\) \term[jingxiang-bianjiezhi]{径向边界值}[radial boundary value]。
\end{theorem}
\symidx{\(f^*\)：Hardy 函数的 \(L^2\) 径向边界值}
\begin{proof}
令 \(F_N(\zeta)=\sum_{n=0}^Na_n\zeta^n\)。圆周正交性给出
\[
 \|F_N-F_M\|_{L^2(m)}^2=\sum_{n=M+1}^N|a_n|^2
 \quad(N>M).
\]
因为 \((a_n)\in\ell^2\)，\((F_N)\) 是 \(L^2\) Cauchy 列。平方可积函数空间的完备性给出极限 \(f^*\)。令 \(N\to\infty\)，有限和的平方范数公式同时给出 \(\|f^*\|_2^2=\sum|a_n|^2\)。这一构造只用非负幂的正交性，尚不需要全部三角系的完备性。

固定 \(r<1\)。\(f_r\) 是 \(\sum_{n\ge0}a_nr^n\zeta^n\) 在圆周上的一致极限，因而也是 \(L^2\) 极限。两列有限和相减并取范数极限，得到\eqref{b3:eq:28-boundary-error}的第一式。对每个 \(n\)，\((1-r^n)^2\to0\)，且它不超过一；可和控制来自 \(\sum|a_n|^2<\infty\)。对计数测度使用支配收敛，便得\eqref{b3:eq:28-hardy-radial-limit}。

范数空间中的极限唯一，所以 \(f^*\) 作为 \(L^2\) 元素唯一。对两个有限和计算内积，再以 Cauchy--Schwarz 控制取极限的误差，得到
\[
 \int_{\mathbb T}f^*\overline{g^*}\,\mathrm dm
 =\sum_{n=0}^{\infty}a_n\overline{b_n}
 =\langle f,g\rangle_{H^2}.
\]
这也表明边界映射是线性等距映射。
\end{proof}

若 \(f\) 连续到闭圆盘，那么 \(f_r\) 一致趋于通常的边界限制 \(f|_{\mathbb T}\)，所以上述 \(f^*\) 就是它的 \(L^2\) 等价类。一般情况下，星号表示通过极限得到的数据，不表示先在圆周上任意给 \(f\) 补值。

\ProseParagraphBreak
需要区分范数极限与逐点径向极限。\eqref{b3:eq:28-hardy-radial-limit}本身只直接保证存在几乎处处收敛的半径子列：选 \(r_j\uparrow1\)，使 \(\|f_{r_j}-f^*\|_2^2\le2^{-j}\)，则 Tonelli 定理给出
\[
 \int_{\mathbb T}\sum_{j=1}^{\infty}|f(r_j\zeta)-f^*(\zeta)|^2
 \,\mathrm dm(\zeta)\le\sum_{j=1}^{\infty}2^{-j}<\infty.
\]
于是误差平方和在几乎处处有限，故 \(f(r_j\zeta)\to f^*(\zeta)\) 几乎处处成立。这里得到沿所选半径子列的几乎处处收敛；所有 \(r\uparrow1\) 的几乎处处径向极限需要进一步的边界定理。

\ProseParagraphBreak
特别地，把有界全纯因子乘进来时，必须说明边界乘积也确实由内部乘积得到。两个 \(L^2\) 极限通常只能直接相乘到 \(L^1\)；这里额外的有界性同时控制了乘积范数与边界值。

\begin{proposition}[有界全纯因子的乘法]\label{b3:prop:28-hardy-multiplier}
设 \(b\in H^\infty(\mathbb D)\)、\(f\in H^2(\mathbb D)\)，并取归一化圆周测度 \(m\)。则 \(bf\in H^2(\mathbb D)\)，且
\[
 \|bf\|_{H^2}\le\|b\|_\infty\cdot\|f\|_{H^2},
 \quad (bf)^*=b^*f^*\quad m\text{-几乎处处}.
\]
此外，\(|b^*|\le\|b\|_\infty\) 几乎处处。若还满足 \(|b^*|=1\) 几乎处处，则
\[
 \|bf\|_{H^2}=\|f\|_{H^2}.
\]
\end{proposition}
\begin{proof}
对每个 \(0<r<1\)，内部的点态界给出
\[
 \int_{\mathbb T}|b(r\zeta)\cdot f(r\zeta)|^2\,\mathrm dm(\zeta)
 \le\|b\|_\infty^2
       \int_{\mathbb T}|f(r\zeta)|^2\,\mathrm dm(\zeta).
\]
取上确界即得 \(bf\in H^2\) 及范数界，且 \(b\in H^2\) 已由有界性保证。对 \(b,f,bf\) 的三份 \(L^2\) 径向误差同时选取半径 \(r_j\uparrow1\)，使它们的平方范数之和小于 \(2^{-j}\)。刚才的可和误差论证说明，除去一个零测集后，三个函数沿这同一子列都收敛到各自的边界函数。在这些圆周点上取通常复数乘法的极限，便有 \((bf)^*=b^*f^*\)。同样，\(|b(r_j\zeta)|\le\|b\|_\infty\) 给出边界模长界。

若 \(|b^*|=1\) 几乎处处，边界等距性给出
\[
 \|bf\|_{H^2}^2
 =\int_{\mathbb T}|b^*|^2\cdot|f^*|^2\,\mathrm dm
 =\int_{\mathbb T}|f^*|^2\,\mathrm dm
 =\|f\|_{H^2}^2.
\]
这一子列上的几乎处处收敛使极限乘积与边界函数相对应。
\end{proof}

例如乘以 \(z\) 只把系数向后移动一位，因而是 \(H^2\) 上的等距映射；边界解释则是 \(|\zeta|=1\)。这个映射不是满射，因为所得函数都在原点为零。范数不变并不意味着不丢失可取到的函数。

\subsection{边界 Fourier 系数}
\label{b3:subsec:280202}

在 \(L^2(\mathbb T,m)\) 中，每个函数也属于 \(L^1\)，因为 \(m(\mathbb T)=1\)。因此对任意整数 \(k\)，其 Fourier 系数
\[
 \widehat F(k)\coloneqq
 \int_{\mathbb T}F(\zeta)\cdot\overline{\zeta^k}\,\mathrm dm(\zeta)
\]
都定义良好。\(\widehat F(k)\) 关于 \(F\) 是连续线性泛函，且 \(|\widehat F(k)|\le\|F\|_2\)。前面的构造表明，解析函数的边界系数只能出现在 \(k\ge0\) 的位置。

\ProseParagraphBreak
这里复用第二卷定理3.3.2与推论3.3.3：在总质量为一的圆周上，\((e^{2\pi\mathrm i kt})_{k\in\mathbb Z}\) 是 \(L^2\) 的标准正交基，Fourier 系数映射是到 \(\ell^2(\mathbb Z)\) 的满射等距映射。令 \(\theta=2\pi t\) 后，基函数变为 \(\zeta^k\)，测度变为 \(\mathrm d\theta/(2\pi)\)，故 Parseval 等式的常数仍为一。这些结论指 \(L^2\) 展开，不要求 Fourier 级数逐点收敛。

\begin{proposition}[边界函数的频率刻画]\label{b3:prop:28-hardy-boundary-range}
取单位圆周 \(\mathbb T\) 及测度 \(\mathrm dm(e^{\mathrm i\theta})=\mathrm d\theta/(2\pi)\)，定义闭子空间
\[
 L^2_+(\mathbb T)\coloneqq
 \{F\in L^2(\mathbb T,m;\mathbb C):\widehat F(k)=0\text{ 对所有 }k<0\}.
\]
边界映射
\[
 H^2(\mathbb D)\longrightarrow L^2_+(\mathbb T),
 \quad f\longmapsto f^*
\]
是满射线性等距映射。若 \(f(z)=\sum_{n\ge0}a_nz^n\)，则
\[
 \widehat{f^*}(k)=
 \begin{cases}a_k,&k\ge0,\\0,&k<0.\end{cases}
\]
任意 \(F\in L^2(\mathbb T,m;\mathbb C)\) 到该子空间的正交投影为
\[
 P_+F=\sum_{n=0}^{\infty}\widehat F(n)\zeta^n
 \quad\text{于 }L^2(\mathbb T,m).
\]
特别地，\(\|P_+F\|_2\le\|F\|_2\)。
\end{proposition}
\symidx{\(L^2_+(\mathbb T)\)：负 Fourier 系数为零的圆周闭子空间}
\symidx{\(P_+\)：圆周 \(L^2\) 到非负频率子空间的正交投影}
\begin{proof}
每个条件 \(\widehat F(k)=0\) 都是一个连续线性泛函的核，因此这些核的交是闭线性子空间。对 \(f^*\) 的多项式逼近取 Fourier 系数，系数泛函的连续性给出所列公式，故边界映射的像包含于 \(L^2_+\)。

反之，设 \(F\in L^2_+\)。由第二卷的 Parseval 等式，\((\widehat F(n))_{n\ge0}\) 平方可和，且
\[
 \sum_{n=0}^{\infty}|\widehat F(n)|^2=\|F\|_2^2.
\]
系数表示定理构造出 \(f(z)=\sum_{n\ge0}\widehat F(n)z^n\in H^2\)。\(f^*\) 与 \(F\) 的所有整数频率系数相同；三角系的完备性说明二者在 \(L^2\) 中相等。于是满射性成立。

对一般 \(F\)，Parseval 等式也保证只保留非负频率的级数在 \(L^2\) 中收敛，所得函数属于 \(L^2_+\)。余项 \(F-P_+F\) 只含负频率，因而与每个非负幂正交，再由非负幂在 \(L^2_+\) 中的稠密性知，它与整个子空间正交。这就是正交投影，且
\[
 \|F\|_2^2=\|P_+F\|_2^2+\|F-P_+F\|_2^2
\]
给出其范数界。
\end{proof}

例如 \(F(\zeta)=\overline\zeta\) 的全部质量在频率 \(-1\)，故 \(P_+F=0\)。它不能是任何圆盘 \(H^2\) 函数的边界值，尽管它光滑且平方可积。又如 \(F(\zeta)=\operatorname{Re}\zeta=(\zeta+\overline\zeta)/2\)，其正交投影为 \(\zeta/2\)，不是原数据本身。边界数据的可积性与它能否由全纯函数产生，是两个不同条件。

\ProseParagraphBreak
还有一个有用的刚性现象：若 \(f\in H^2\) 的边界值实值几乎处处，则 \(f\) 必为实常数。因为实值 \(F\) 满足 \(\widehat F(-n)=\overline{\widehat F(n)}\)，负频率消失就迫使全部非零频率消失。这里结论依赖边界函数来自 \(H^2\)；仅给某个全纯函数零测集以外的边界行为，不能不加增长条件就代替这项假设。

\subsection{Cauchy 表示与点值控制}
\label{b3:subsec:280203}

经典 Cauchy 公式在内部圆周上成立，那里边值连续。本节的边界函数只有平方可积性。要把圆周移到 \(\mathbb T\)，应先检验积分核的可积性与极限，而不是把经典公式中的半径直接换成一。

\begin{proposition}[边界 Cauchy 表示与最优点值界]\label{b3:prop:28-hardy-cauchy}
设 \(f\in H^2(\mathbb D)\)，\(f^*\) 是其相对于归一化圆周测度 \(m\) 的 \(L^2\) 边界值。对任意 \(w\in\mathbb D\)，有
\begin{equation}\label{b3:eq:28-hardy-cauchy}
 f(w)=\int_{\mathbb T}
       \frac{f^*(\zeta)}{1-w\overline\zeta}\,\mathrm dm(\zeta),
\end{equation}
且
\begin{equation}\label{b3:eq:28-hardy-point-evaluation}
 |f(w)|\le\frac{\|f\|_{H^2}}{\sqrt{1-|w|^2}}.
\end{equation}
对每个固定 \(w\)，右端常数是最优的；除零函数外，等号成立当且仅当
\[
 f(z)=\frac{c}{1-\overline w z}
 \quad(z\in\mathbb D),\quad c\in\mathbb C\setminus\{0\}.
\]
\end{proposition}
\begin{proof}
固定 \(w\in\mathbb D\)。由于 \(|1-w\overline\zeta|\ge1-|w|>0\)，积分核有界，而 \(f^*\in L^1(m)\)，所以积分绝对存在。几何级数
\[
 \frac1{1-w\overline\zeta}
 =\sum_{n=0}^{\infty}w^n\overline\zeta^{\,n}
\]
在 \(\zeta\in\mathbb T\) 上绝对一致收敛。乘以 \(f^*\) 后，用可积控制 \(|f^*|/(1-|w|)\) 交换积分与求和，即得
\[
 \int_{\mathbb T}\frac{f^*(\zeta)}{1-w\overline\zeta}\,\mathrm dm
 =\sum_{n=0}^{\infty}w^n\widehat{f^*}(n)
 =\sum_{n=0}^{\infty}a_nw^n=f(w).
\]
这个证明同时解释了为什么表示式只读取非负 Fourier 系数。

函数 \(K_w(\zeta)=1/(1-\overline w\zeta)\) 的圆周平方范数由几何级数和正交性给出：
\[
 \|K_w\|_{L^2(m)}^2
 =\sum_{n=0}^{\infty}|w|^{2n}
 =\frac1{1-|w|^2}.
\]
于是\eqref{b3:eq:28-hardy-cauchy}是 \(\langle f^*,K_w\rangle_{L^2}\)，Cauchy--Schwarz 不等式与边界等距性给出点值界。对非零 \(f\)，等号成立当且仅当 \(f^*=cK_w\)；边界映射的单射性使这等价于 \(f(z)=c/(1-\overline w z)\)。这些函数确实属于 \(H^2\)，所以常数不能减小。
\end{proof}

若 \(f\) 连续到闭圆盘，令 \(\zeta=e^{\mathrm i\theta}\) 并利用 \(\mathrm d\zeta=\mathrm i\zeta\,\mathrm d\theta\)，就会把\eqref{b3:eq:28-hardy-cauchy}写成通常的 \((2\pi\mathrm i)^{-1}\int_{\mathbb T}f(\zeta)/(\zeta-w)\,\mathrm d\zeta\)。对于一般 \(f\in H^2\)，Lebesgue 积分中的边界函数是一个 \(L^2\) 等价类，改变零测集上的值不会改变积分。

\ProseParagraphBreak
同一个边界函数还有 Poisson 表示。以本卷\secref{b3:sec:1601}的单位圆盘归一化记核
\[
 P(w,\zeta)=\frac{1-|w|^2}{|\zeta-w|^2}
 =1+\sum_{n\ge1}\bigl(w^n\overline\zeta^{\,n}
                     +\overline w^{\,n}\zeta^n\bigr).
\]
对固定 \(w\in\mathbb D\)，该级数在圆周上一致绝对收敛，故它也可与 \(f^*\in L^1\) 逐项积分。含 \(\overline w^{\,n}\) 的部分乘以 \(\widehat{f^*}(-n)=0\)，从而
\begin{equation}\label{b3:eq:28-hardy-poisson}
 f(w)=\int_{\mathbb T}P(w,\zeta)\cdot f^*(\zeta)\,\mathrm dm(\zeta).
\end{equation}
这只是把本处 \(L^2\) 数据代入一个有界内部核，并没有引用前章只针对连续数据的边界收敛结论。

\ProseParagraphBreak
对于任意 \(F\in L^2\)，Cauchy 积分读取 \(P_+F\)，Poisson 积分则保留正、负两组频率。以 \(F(\zeta)=\overline\zeta\) 为例，前者为零，后者为 \(\overline w\)。只有当负频率消失时，两种内部表示才相同。因此 Poisson 延拓通常是调和的，全纯性还需要边界数据满足频率条件。

\subsection{Hardy 核与再生性质}
\label{b3:subsec:280204}

点值估计说明，在 \(H^2\) 中取某个内点的函数值是一个连续线性操作。Cauchy 表示里的核不仅属于边界的 \(L^2\) 空间，它本身也是一个圆盘全纯函数。由此可以把点值直接表示成两个解析函数的内积。

\begin{definition}[再生核]\label{b3:def:28-reproducing-kernel}
设 \(X\) 是非空集合，\(\mathcal H\) 是由 \(X\) 上复值函数组成的复 Hilbert 空间，内积对第一变量线性。若函数 \(K:X\times X\to\mathbb C\) 满足：对每个 \(w\in X\)，函数 \(K_w=K(\cdot,w)\) 属于 \(\mathcal H\)，且
\[
 f(w)=\langle f,K_w\rangle_{\mathcal H}
 \quad(f\in\mathcal H),
\]
则称 \(K\) 为 \(\mathcal H\) 的\term[zaisheng-he]{再生核}[reproducing kernel]，并称 \(\mathcal H\) 为\term[zaisheng-he-hilbert-kongjian]{再生核 Hilbert 空间}[reproducing kernel Hilbert space]。
\end{definition}

定义中的元素是实际的函数，因而每个点值已经有意义。一般 \(L^2\) 空间的元素只按几乎处处相等区分，单点值不能由其等价类确定；不能把任意 \(L^2\) 空间直接当作再生核空间。本章的 \(H^2\) 元素始终是在盘内逐点定义的全纯函数，边界映射才把它们识别为一个特殊的 \(L^2\) 子空间。

\begin{theorem}[Hardy 核]\label{b3:thm:28-hardy-kernel}
在单位圆盘 \(\mathbb D\) 上，给 \(H^2(\mathbb D)\) 配备内积
\(\langle f,g\rangle=\int_{\mathbb T}f^*\overline{g^*}\,\mathrm dm\)，其中 \(\mathrm dm(e^{\mathrm i\theta})=\mathrm d\theta/(2\pi)\)。其再生核为
\begin{equation}\label{b3:eq:28-hardy-kernel}
 S(z,w)\coloneqq\frac1{1-z\overline w},
 \quad z,w\in\mathbb D.
\end{equation}
对每个 \(w\in\mathbb D\)，记 \(k_w=S(\cdot,w)\)。则对任意 \(f\in H^2(\mathbb D)\)，有
\[
 f(w)=\langle f,k_w\rangle_{H^2},
 \quad \|k_w\|_{H^2}^2=S(w,w)=\frac1{1-|w|^2}.
\]
该核称为圆盘的\term[hardy-he]{Hardy 核}[Hardy kernel]，也称\term[szego-he]{Szegő 核}[Szegő kernel]。
\end{theorem}
\symidx{\(S(z,w)\)：单位圆盘的 Hardy 核}
\begin{proof}
固定 \(w\in\mathbb D\)。\(k_w(z)=\sum_{n\ge0}\overline w^{\,n}z^n\) 的系数平方和为 \((1-|w|^2)^{-1}\)，所以它属于 \(H^2\)。若 \(f(z)=\sum a_nz^n\)，则
\[
 \langle f,k_w\rangle_{H^2}
 =\sum_{n=0}^{\infty}a_n\overline{\overline w^{\,n}}
 =\sum_{n=0}^{\infty}a_nw^n=f(w).
\]
再取 \(f=k_w\) 就得到 \(\|k_w\|^2=k_w(w)=S(w,w)\)。

再生性还保证唯一性：若 \(\widetilde k_w\) 也代表同一点值，则 \(\langle f,k_w-\widetilde k_w\rangle=0\) 对所有 \(f\in H^2\) 成立。取 \(f=k_w-\widetilde k_w\)，可得两函数的距离为零，故相等。
\end{proof}

核在 \(z\) 上全纯，在 \(w\) 上反全纯；\(\overline w\) 的位置由内积的线性约定决定。交换两个变量有 \(S(z,w)=\overline{S(w,z)}\)，而
\[
 \langle k_w,k_z\rangle_{H^2}=k_w(z)=S(z,w).
\]
这些等式也说明，核的两个变量并非两个独立复制的积分变量：第一个变量表示函数的自变量，第二个变量指定要恢复的点值。

\ProseParagraphBreak
核函数还给出一个可直接求解的极值问题。给定 \(w\in\mathbb D\) 与 \(c\in\mathbb C\)，在满足 \(f(w)=c\) 的全部 \(f\in H^2\) 中，唯一范数最小的函数是
\[
 f_0(z)=c(1-|w|^2)\cdot k_w(z),
 \quad \|f_0\|_{H^2}=|c|\sqrt{1-|w|^2}.
\]
确实，\(f_0(w)=c\)；对任一同样满足该约束的 \(f\)，有 \(\langle f-f_0,k_w\rangle=(f-f_0)(w)=0\)，因此
\[
 \|f\|_{H^2}^2=\|f_0\|_{H^2}^2+\|f-f_0\|_{H^2}^2.
\]
最小性与唯一性同时得出。再生核把一个点值约束变成 Hilbert 空间中的正交分解；\secref{b3:sec:2805}会用面积范数得到另一种核，并比较同一点值的代价如何改变。
